\documentclass[12pt,letterpaper]{article}

% ---------------- PAQUETES ----------------
\usepackage[spanish,es-tabla]{babel}
\usepackage[utf8]{inputenc}
\usepackage[T1]{fontenc}

\usepackage[left=2.5cm,right=2.5cm,top=2.5cm,bottom=2.5cm]{geometry}
\usepackage{graphicx}
\usepackage{fancyhdr}
\usepackage{titlesec}
\usepackage{xcolor}
\usepackage{booktabs}
\usepackage{array}
\usepackage{parskip}
\usepackage{enumitem}
\usepackage{amssymb}
\usepackage{hyperref}


% ---------------- FORMATOS ----------------
\titleformat{\section}{\large\bfseries}{\thesection.}{1em}{}
\titleformat{\subsection}{\normalsize\bfseries}{\thesubsection.}{1em}{}

% ---------------- DOCUMENTO ----------------
\begin{document}
\section{Introducción}
El presente documento tiene como objetivo identificar las principales problemáticas dentro del proceso de elección de subdirección en la universidad, así como proponer soluciones tecnológicas que permitan mejorar la participación, transparencia y accesibilidad del proceso electoral.

% ---------------- PROBLEMATICAS ----------------
\section{Problemáticas Detectadas}

\begin{table}[h!]
\centering
\begin{tabular}{p{0.8cm} p{11cm}}
\toprule
\textbf{No.} & \textbf{Problemática} \\
\midrule
1 & Desconocimiento del proceso electoral \\
2 & Falta de difusión \\
3 & Proceso poco simplificado \\
4 & Baja motivación del alumnado \\
5 & Escasa simpatía con el alumnado \\
6 & Poco acceso a las propuestas \\
7 & Dudas sobre la seguridad del proceso \\
8 & Falta de integración de la comunidad \\
\bottomrule
\end{tabular}
\end{table}

% ---------------- ANALISIS ----------------
\section{Análisis de Causas}

\begin{table}[h!]
\centering
\begin{tabular}{p{0.8cm} p{5.5cm} p{5.5cm}}
\toprule
\textbf{No.} & \textbf{Problemática} & \textbf{Causa Principal} \\
\midrule
1 & Desconocimiento del proceso & Información poco accesible \\
2 & Falta de difusión & Falta de interés institucional y estudiantil \\
3 & Proceso poco simplificado & Falta de claridad y objetividad \\
4 & Baja motivación & Ausencia de incentivos o consecuencias \\
5 & Escasa simpatía & Falta de conexión con la comunidad \\
6 & Poco acceso a propuestas & Canales de comunicación limitados \\
7 & Seguridad & Necesidad de garantizar anonimato \\
8 & Falta de integración & Escasa participación activa \\
\bottomrule
\end{tabular}
\end{table}

% ---------------- SOLUCIONES ----------------
\section{Propuesta de Soluciones Tecnológicas}

Se propone el desarrollo de una aplicación web/móvil que centralice todo el proceso electoral, facilitando la participación estudiantil y mejorando la transparencia.

\begin{table}[h!]
\centering
\begin{tabular}{p{0.8cm} p{5cm} p{6cm}}
\toprule
\textbf{No.} & \textbf{Funcionalidad} & \textbf{Descripción} \\
\midrule
1 & Portal informativo & Sección clara que explique el proceso paso a paso con interfaz intuitiva \\
2 & Notificaciones & Envío de avisos sobre fechas importantes mediante correo o app \\
3 & Perfil de candidatos & Visualización de propuestas, trayectoria y videos explicativos \\
4 & Comparador de propuestas & Herramienta para analizar diferencias entre candidatos \\
5 & Sistema de votación digital & Plataforma segura con autenticación de usuarios \\
6 & Anonimato garantizado & Separación entre identidad del usuario y voto emitido \\
7 & Métricas en tiempo real & Visualización de participación (sin comprometer privacidad) \\
8 & Espacio de interacción & Foros o sección de preguntas para candidatos \\
\bottomrule
\end{tabular}
\end{table}

% ---------------- VALOR ----------------
\section{Valor de la Propuesta}
El desarrollo de esta solución tecnológica permitiría:
\begin{itemize}
    \item Incrementar la participación estudiantil
    \item Mejorar la transparencia del proceso
    \item Facilitar el acceso a la información
    \item Generar mayor confianza en los resultados
\end{itemize}

% ---------------- CONCLUSION ----------------
\section{Conclusión}
La implementación de herramientas tecnológicas dentro del proceso electoral universitario representa una oportunidad para modernizar la participación estudiantil. Una aplicación bien diseñada puede resolver múltiples problemáticas actuales y fomentar una comunidad más informada, activa y comprometida.

\end{document}